\begin{enumerate}

\item Formar la matriz del tensor de Cauchy y obtener las tensiones principales en el plano xy.

\item Añadir la tensión normal sigma z, ordenar los tres autovalores y calcular la máxima cortadura absoluta.

\item Obtener la tensión equivalente de von Mises y la magnitud equivalente de Tresca.

\item Contrastar los resultados con los invariantes del tensor.

\item Si se proporciona el límite elástico, comparar ambos criterios de fluencia.

\end{enumerate}

\noindent\textbf{Operaciones y resultados intermedios.}

\begin{enumerate}

\item Centro del círculo en el plano xy.

\[\fitmath{y_{1}=\frac{-90+25}{2}\approx -32.5\,\mathrm{MPa}}\]

\item Radio del círculo de Mohr.

\[\fitmath{y_{2}=\sqrt{{\left(\frac{-90-25}{2}\right)}^{2}+{\left(-55\right)}^{2}}\approx 79.56915\,\mathrm{MPa}}\]

\item Tensión principal mayor en el plano.

\[\fitmath{y_{3}=\frac{-90+25}{2}+\sqrt{{\left(\frac{-90-25}{2}\right)}^{2}+{\left(-55\right)}^{2}}\approx 47.06915\,\mathrm{MPa}}\]

\item Tensión principal menor en el plano.

\[\fitmath{y_{4}=\frac{-90+25}{2}-\sqrt{{\left(\frac{-90-25}{2}\right)}^{2}+{\left(-55\right)}^{2}}\approx -112.0692\,\mathrm{MPa}}\]

\item Von Mises a partir de las componentes.

\[\fitmath{y_{5}=\frac{\sqrt{{\left(-90-25\right)}^{2}+{\left(25--40\right)}^{2}+{\left(-40--90\right)}^{2}+6\cdot {\left(-55\right)}^{2}}}{\sqrt{2}}\approx 138.0217\,\mathrm{MPa}}\]

\item Índice tensional de Tresca.

\[\fitmath{y_{6}=2\cdot \frac{5\cdot \sqrt{1013}}{2}\approx 159.1383\,\mathrm{MPa}}\]

\item Coeficiente de seguridad de von Mises.

\[\fitmath{y_{7}=\frac{260}{\frac{\sqrt{2}\cdot \sqrt{{\left(-15+5\cdot \sqrt{1013}\right)}^{2}+{\left(15+5\cdot \sqrt{1013}\right)}^{2}+101300}}{4}}\approx 1.883761\,}\]

\item Coeficiente de seguridad de Tresca.

\[\fitmath{y_{8}=\frac{260}{5\cdot \sqrt{1013}}\approx 1.633799\,}\]

\item Respuesta principal: comprobación final.

\[\fitmath{y_{9}=\frac{\frac{-65}{2}+\frac{5\cdot \sqrt{1013}}{2}-\left(\frac{-5\cdot \sqrt{1013}}{2}-\frac{65}{2}\right)}{2}\approx 79.56915\,\mathrm{MPa}}\]

\end{enumerate}
